\section{Introduction}
\IEEEPARstart{H}{AND} exoskeletons are robotic assistance and rehabilitation
devices whose function is to accompany or help recover the movement of the
fingers \cite{ref_exo_review}. For the assistance to be comfortable and safe,
the mechanism must faithfully reproduce the natural joint trajectories of the
finger over the whole flexion range; otherwise, unwanted reaction forces appear
at the joints and the device stops being comfortable for the user
\cite{ref_fourbar_nat}. Achieving that tracking depends largely on correctly
choosing the link lengths of the mechanism, a problem known as
\emph{dimensional synthesis} \cite{ref_de_fourbar}.

Dimensional synthesis for path generation is a hard optimization problem: the
solution space contains large regions where the mechanism simply does not
assemble or reverses its motion, and the relation between the link lengths and
the curve described by the mechanism is nonlinear and has many local optima. For
this reason, classical gradient-based methods tend to get trapped in suboptimal
solutions, which motivates the use of population-based metaheuristics such as
Differential Evolution (DE) and Genetic Algorithms (GA), which can traverse the
search space without needing derivative information
\cite{ref_de_fourbar,ref_tlbo}.

Although both optimization techniques are widely used, their effectiveness
depends on the nature of the problem, on how the variables are represented and on
the tuning of the hyperparameters \cite{ref_ga_pso_de,ref_de_vs_ga}. The choice
between one or the other is rarely obvious and almost always requires testing
them on the specific case to be solved.

Dimensional synthesis of linkages for path generation has historically been
addressed with analytical methods and, more recently, with metaheuristics that
avoid the limitations of gradient-based approaches on nonconvex and discontinuous
functions \cite{ref_de_fourbar,ref_tlbo}. Four-bar mechanisms are a recurring
building block both in path synthesis and in finger exoskeletons, due to their
simplicity and their ability to approximate complex coupler curves with few
parameters \cite{ref_fourbar_nat,ref_exo_stephenson}. When several four- and
five-bar stages are chained, the number of variables grows and the search space
becomes strongly multimodal, which makes the use of global optimization
techniques unavoidable.

In the field of hand exoskeletons, several recent works use swarm-based or
evolutionary optimization to tune the dimensions of mechanisms that reproduce the
natural flexion of the fingers \cite{ref_exo_review,ref_exo_stephenson}. The
choice of algorithm, however, is rarely justified empirically on the specific
mechanism. The direct comparison between families of algorithms has also been
studied: in continuous-domain problems, DE and particle swarm optimization tend
to behave more naturally than the GA, which was originally conceived for discrete
representations \cite{ref_ga_pso_de,ref_de_vs_ga}. Nevertheless, the outcome of
the comparison depends on the encoding, on the operators and on the tuning of
hyperparameters, so it remains necessary to characterize the behavior on real
instances \cite{ref_de_survey,ref_de_sade}.

Unlike those works, here the same formulation, objective function and number of
generations are fixed for both algorithms on a real exoskeleton mechanism, and
not only the fitting error is evaluated but also the coherence and compactness of
the resulting mechanism, two aspects that are decisive for its later manufacture.

In this work both algorithms are used under equal conditions, that is, starting
from the same formulation, the same objective function and the same number of
generations, so that any difference in the result is due to the algorithm and not
to how the experiment was set up. This with the aim of determining the best
technique to synthesize planar mechanisms.

For the comparison to be reproducible, a \emph{partial} version of the mechanism
is optimized, which generates the movement up to the medial phalanx (PIP joint).
This case is smaller than the full three-phalanx model, but it keeps the
essential difficulties of the problem (the coupling between stages, the assembly
constraints and the need for the movement not to reverse), which makes it
possible to study the effect of the optimization algorithm without the added
complexity of the full model.

\subsection{Contributions}
The contributions of this article are: (i) a formulation of the partial
dimensional synthesis problem with 16 parameters, an objective function based on
the weighted Chamfer distance and feasibility penalties; (ii) a controlled
comparison of DE and GA under identical conditions, which evaluates not only the
fitting error but also the coherence and compactness of the resulting mechanism;
and (iii) an analysis of the reasons why DE keeps a slight advantage in this
continuous-variable problem.
